\documentclass{article}
\usepackage{amsmath, amssymb}
\usepackage{array}
\usepackage{xcolor}

\begin{document}

\section{Introduction}

Lecture 5 is about relative entropy and differential entropy.

\section{Relative entropy}

It's really just a generalization of entropy.

\subsection{Definition}

Let $p=\left(p_1,\ldots,p_m\right)$ and $q=\left(q_1,\ldots, q_m\right)$ are two distributions
\[
\operatorname{D}\left(p\| q\right) = \sum_{i} p_i\log\left(\frac{p_i}{q_i}\right) \geq 0
\]
Note that 
\begin{itemize}
    \item $0\log \frac{0 }{0 }= 0$
    \item $a\log \frac{a }{0 }=\infty$ (for $a>0$)
    \item $\operatorname{D}\left(p\| q\right)$ may be infinite
\end{itemize}
There are many names and many different interpretations.

\subsection{Intuition}

Let $X$ be a random variable over $\left\{1,\ldots, m\right\}$. We know that $\operatorname{H}(X)$ measures the uncertainty of $X$, the amount of information that we don't know about $X$. We can view $\log m - \operatorname{H}(X)$ as the measurement of the information that we do have about $X$, simply by knowing the distribution.

\subsubsection{Example}

Let $m=4$ and $X=\left(X_1, X_2\right)\in \left\{00,01,10,11\right\}$. Let the distribution be $\left(\frac{1}{2}, 0, 0, \frac{1 }{2}\right)$. It follows that ... and I missed it.

\subsubsection{Back to interpretation}

Originally we thought that $X$ is uniformly distributed. We learned that $X\sim \left(p_1,\ldots,p_m\right)$. How mu information did we gain about $X$? 
\[
\operatorname{H}\left(\frac{1}{m},\ldots,\frac{1 }{m}\right) - \operatorname{H}\left(p_1,\ldots,p_m\right)=\log m -\operatorname{H}(X)
\]
Note that 
\[
\log m -\operatorname{H}\left(p_1,\ldots,p_m\right)= \log m + \sum_i p_i \log p_i = \sum_i p_i\log \left(\frac{p_i }{1/m}\right)=\operatorname{D}\left(p\| q\right)
\]
where $q$ is the uniform distribution.

Thus $\operatorname{D}\left(p\| q\right)$, where $q$ is the uniform distribution, is a measure for the information that we gained about $X$ if originally we thought that $X$ is uniformly distributed and then learned what the true distribution is.

\subsection{Interpretation of $\operatorname{D}\left(p\| q\right)$}

More generally: Assume that we originally thought that a random variable $X$ is distributed according to $q$, and we learned that it is distributed according to $p$. The relative entropy $\operatorname{D}\left(p\| q\right)$ is a measure for the amount of information that we gained.

\subsection{Today's Lecture}

We will try to better understand this notion, its relevance, and why it is used to measure the information gain.
\begin{itemize}
    \item What does it mean that we learned that $X$ is distributed differently? How can the distribution of $X$ change?
    \item Why this function?
    \item What are its properties?
    \item Are there interesting relations to mutual information?
\end{itemize}

\section{How can the distribution change?}

What does it mean that originally we thought that $X$ is distributed according to $q$, and then we learned that it is distributed according to $p$? How can the distribution of $X$ change?

\textbf{Usually, this happens when we learn additional information about $X$.} We condition on something.

\subsection{Example}

Let $X\in\left\{1,2,3,4\right\}$. The distribution of $X$ is $q=\left(\frac{1 }{2}, \frac{1 }{4 },\frac{1 }{4 },0\right)$ ... and I missed the rest.

\subsubsection{Another example}

Missed it too, see photo.

\subsection{$\operatorname{D}\left(p\| q\right)$ versus $\operatorname{H}(q)-\operatorname{H}(p)$}

We saw that if $q$ is the uniform distribution,
\[
\operatorname{D}\left(p\| q\right) = \log m - \operatorname{H}(p) = \operatorname{H}(q)-\operatorname{H}(p)
\]
... and I missed the rest, see photo.

\subsection{Example}

If $q$ is the distribution of $X$, and $E$ is an event that depends only on $X$, and $p$ is the distribution of $X$ conditioned on $E$. That is $q=P_X$ and $p=P_{X\mid E}$
\[
p_i = \begin{cases}
   0 & i\notin E \\
   \frac{q_i}{P_X(E)} & i\in E
\end{cases}
\]
Then 
\[
\operatorname{D}\left(p\| q\right) = \operatorname{D}\left(P_{X\mid E}\| P_X\right) = \log \frac{1}{P_X(E)}
\]
is the information that we gained by learning that $X\in E$.
\begin{align*}
    \operatorname{D}\left(p\| q\right) &= \sum_i p_i\log \left(\frac{p_i }{q_i}\right) = \sum_{i\in E} p_i\log\left(\frac{p_i }{q_i}\right)\\
    &= \sum_{i\in E}p_i  \log \frac{1}{P_X(E)} =  \log \frac{1}{P_X(E)}
\end{align*}
I'm unsure what we just did here.

\section{Does $\operatorname{D}\left(p\| q\right)=\infty$ make sense?}

If for some $i$ we have $q_i=0$ and $p_i>0$ then $\operatorname{D}\left(p\| q\right)=\infty$. Makes sense because if originally the probability for $i$ is $0$, it cannot become positive by learning something (it cannot become positive by conditioning on an event).

Hence, it makes sense to think of the information gain as infinite: We were absolutely sure that $X\neq i$ and now we think that maybe $X=i$.

Alternatively, we can define $\operatorname{D}\left(p\| q\right)$ only for distributions such that $q_i=0$ implies $p_i=0$.

\section{Axiomatic Derivation: The Axioms}

Why this function? Well let's do an axiomatic derivation.
\begin{enumerate}
    \item $\operatorname{D}\left(p\| q\right)$ is continuous and symmetric (invariant to permutations of coordinates, not $p$ and $q$)
    \item If $q=\left(\frac{1}{m},\ldots,\frac{1}{m}\right)$, then $\operatorname{D}\left(p\| q\right)=\log m -\operatorname{H}\left(p_1,\ldots,p_m\right)$
    \item \textbf{Special case} \begin{align*}
\operatorname{D}\bigl( (p_1,\ldots,p_m) &\parallel (q_1,\ldots,q_m) \bigr) = \\
\operatorname{D}\biggl( \Big(p_1,\ldots,p_{m-1},\frac{p_m}{2},\frac{p_m}{2}\Big) &\parallel \Big(q_1,\ldots,q_{m-1},\frac{q_m}{2},\frac{q_m}{2}\Big) \biggr)
\end{align*}
\end{enumerate}
Intuitively, symbol $m$ is partitioned into two new symbols with equal probabilities in both $p,q$. We gained new information about about the original partition to $\left\{1,\ldots, m\right\}$ but no new information about the partition of symbol $m$ into two symbols.

\subsection{General case}
\begin{align*}
\operatorname{D}\bigl( (p_1,\ldots,p_m) &\parallel (q_1,\ldots,q_m) \bigr) = \\
\operatorname{D}\biggl( \Big(\frac{p_1}{k_1},\ldots,\frac{p_1}{k_1},\ldots,\frac{p_m}{k_m},\ldots\frac{p_m}{k_m}\Big) &\parallel \Big(\frac{q_1}{k_1},\ldots,\frac{q_1}{k_1},\ldots,\frac{q_m}{k_m},\ldots\frac{q_m}{k_m}\Big) \biggr)
\end{align*}
where each $\frac{p_i}{k_i}$ and $\frac{q_i}{k_i}$ are repeated $k_i$ times.

Intuitively, we only gained new information about the partition into groups but not about the more refined partition inside each group. It's kind of just the grouping axiom, but for relative entropy.

This can also be viewed as a special case of the chain rule. $p,q$ in the right hand side can be viewed as joint distribution of $X,Y$, where $X$ indicates the group and $Y$ indicates the place in the group. Only the distribution of $X$ changed (and not the distribution of $Y$ conditioned on $X$), so the information gained is only from the change in the distribution of $X$.

\subsection{Axiomatic Derivation of $\operatorname{D}\left(p\| q\right)$}

Let $\left(p_1,\ldots,p_m\right)$ be an arbitrary distribution. Let $\left(q_1,\ldots, q_m\right)$ be an arbitrary distribution, such that $q_1,\ldots,q_m$ are positive and rational, $q_i=\frac{k_i }{M}$ for some $k_1,\ldots,k_m,M>0$ (where $M$ is a common multiple).

Assume $q_i = \frac{k_i }{M }>0$. By the third axiom, 
\begin{align*}
\operatorname{D}\bigl( (p_1,\ldots,p_m) &\parallel (q_1,\ldots,q_m) \bigr) = \\
\operatorname{D}\biggl( \Big(\frac{p_1}{k_1},\ldots,\frac{p_1}{k_1},\ldots,\frac{p_m}{k_m},\ldots\frac{p_m}{k_m}\Big) &\parallel \Big(\frac{q_1}{k_1},\ldots,\frac{q_1}{k_1},\ldots,\frac{q_m}{k_m},\ldots\frac{q_m}{k_m}\Big) \biggr)
\end{align*}
where each $\frac{p_i}{k_i}$ and $\frac{q_i}{k_i}$ are repeated $k_i$ times. It follows that 
\begin{align*}
\operatorname{D}\bigl( (p_1,\ldots,p_m) &\parallel (q_1,\ldots,q_m) \bigr) = \\
\operatorname{D}\biggl( \Big(\frac{p_1}{k_1},\ldots,\frac{p_1}{k_1},\ldots,\frac{p_m}{k_m},\ldots\frac{p_m}{k_m}\Big) &\parallel \Big(\frac{1}{M},\ldots,\frac{1}{M},\ldots,\frac{1}{M},\ldots\frac{1}{M}\Big) \biggr)
\end{align*}
By the second axiom, this is equal to 
\[
\log M - \operatorname{H}\left(\frac{p_1}{k_1},\ldots,\frac{p_1}{k_1},\ldots,\frac{p_m}{k_m},\ldots\frac{p_m}{k_m}\right)
\]
Assume $q_i = \frac{k_i }{M }>0$. Thus,
\begin{align*}
    \operatorname{D}\left(p\| q\right) &= \log M - \operatorname{H}\left(\frac{p_1}{k_1},\ldots,\frac{p_1}{k_1},\ldots,\frac{p_m}{k_m},\ldots\frac{p_m}{k_m}\right)\\
    &= \sum_{i=1 }^{m }p_i\log M + \sum_{i=1 }^{m }p_i\log\left(\frac{p_i }{k_i }\right) = \sum_{i=1 }^{m }p_i\log\left(\frac{p_i }{k_i/M }\right)\\
    &= \sum_{i=1 }^{m }p_i\log \left(\frac{p_i }{q_i}\right)
\end{align*}
We assumed that for every $i,q_i>0 $ is rational. The irrational case and the case $q_i=0$ follow by continuity.

\section{The information that $Y$ gives on $X$}

The information that the event $Y=b$ gives on $X$ is 
\[
\operatorname{D}\left(P_{X\mid Y=b}\| P_X\right)
\]
Averaging over all elements $b$, we get the information that $Y$ gives on $X$
\begin{align*}
    \sum_{b} \operatorname{P}\left(Y=b\right)\cdot & \operatorname{D}\left(P_{X\mid Y=b}\| P_X\right)\\
     &= \sum_b \operatorname{P}\left(Y=b\right)\sum_a \operatorname{P}\left(X=a\mid Y=b\right)\log \frac{\operatorname{P}\left(X=a\mid Y=b\right)}{\operatorname{P}\left(X=a\right)}\\
    &= \sum_{a,b} \operatorname{P}\left(X=a, Y=b\right)\log \frac{\operatorname{P}\left(X=a\mid Y=b\right)}{\operatorname{P}\left(X=a\right)}\\
    &= \operatorname{H}(X) -= \operatorname{H}(X\mid Y)=\operatorname{I}(X;Y)
\end{align*}

\subsection{An Additional Connection to Mutual Information}

\[
\operatorname{I}(X;Y) = \operatorname{D}\left(P_{X,Y}\| P_X \times P_Y\right)
\]
that is, the entropy of $\operatorname{P}\left(X=a,Y=b\right)$, relative to the product distribution $\operatorname{P}\left(X=a\right)\cdot \operatorname{P}\left(Y=b\right)$ (as if $Y,X$ were independent).

\textbf{The proof follows immediately by the definitions.}

\section{Conditional Relative Entropy}

$P,Q$: two joint probability distributions of two random variables $X,Y$ (could be thought of as two matrices of values) (a good example to have in mind is that $Q$ is the original probability distribution and $P$ is the distribution $Q$ conditioned on an event $E$)

For every event $Y=b$, we can consider the relative entropy of the distributions $P_{X\mid Y=b}$ and $Q_{X\mid Y=b}$
\[
\operatorname{D}\left(P_{X\mid Y=b}\| Q_{X\mid Y=b}\right)=\sum_a P(X=a\mid Y=b)\log \frac{P(X=a\mid Y=b)}{Q(X=a\mid Y=b)}
\]
The conditional relative entropy of the distributions $P_{X\mid Y}$ and $Q_{X\mid Y}$ is defined by averaging over all elements $b$,
\begin{align*}
    \operatorname{D}\left(P_{X|Y}\| Q_{X|Y}\right) &= \sum_b P(Y=b)\cdot \operatorname{D}\left(P_{X\mid Y=b}\| Q_{X\mid Y=b}\right)\\
    &= \sum_{a,b} P(X=a, Y=b)\log \frac{P(X=a\mid Y=b)}{Q(X=a\mid Y=b)}
\end{align*}
Sometimes denoted as 
\[
\operatorname{D}\left(P(X\mid Y)\| Q(X\mid Y)\right)
\]

\subsection{Chain Rule}

For two joint probability distributions $P,Q$ of two random variables $X,Y$
\[
\operatorname{D}\left(P_{X,Y}\| Q_{X,Y}\right) = \operatorname{D}\left(P_Y\| Q_Y\right) + \operatorname{D}\left(P_{X|Y}\| Q_{X|Y}\right)
\]
The chain rule generalizes to $n$ variables $X_1,\ldots, X_n$.

\subsection{Superadditivity}

For two joint probability distributions $P,Q$ of two random variables $X,Y$ such that $Q_{X,Y}$ is a product distribution: $Q_{X,Y} = Q_X \times Q_Y$ (that is $X,Y$ are independent according to $Q$)
\[
\operatorname{D}\left(P_{X,Y}\| Q_{X,Y}\right) \geq \operatorname{D}\left(P_{X}\| Q_{X}\right)  + \operatorname{D}\left(P_{Y}\| Q_{Y}\right) 
\]
The proof follows by a convexity argument. The direction of the inequality is reversed, compared to subadditivity of entropy.

Superadditivity is only true for a product distribution $Q$. In general, unlike entropy $\operatorname{D}\left(P_{X|Y}\| Q_{X|Y}\right)$ and $\operatorname{D}\left(P_{X}\| Q_{X}\right)$ are incomparable.

\subsubsection{Generalizes to $n$ random variables}

Superadditivity generalizes to $n$ variables $X_1,\ldots, X_n$: For two joint probability distributions $P,Q$ of $n$ random variables $X_1,\ldots, X_n$ such that $Q_{X_1,\ldots,X_n}$ is a product distribution:
\[
Q_{X_1,\ldots,X_n} = Q_{X_1}\times \cdots \times Q_{X_n}
\]
that is $X_1,\ldots, X_n$ are independent according to $Q$. We have ... and I missed it.

\section{Convexity}

For four probability distributions $p_1,p_2,q_1,q_2$, over the same finite set, and every $0\leq \lambda \leq 1$, let 
\[
p=\lambda p_1 + (1-\lambda) p_2
\]
and 
\[
q= \lambda q_1 + (1-\lambda) q_2
\]
Then,
\[
\operatorname{D}\left(p\| q\right) \leq \lambda \operatorname{D}\left(p_1\| q_1\right)+ (1-\lambda)\operatorname{D}\left(p_2\| q_2\right)
\]

\section{Relative Entropy as a Distance}

$\operatorname{D}\left(p\| q\right)$ is sometimes called KL-distance and used as a measure for how far $p$ is from $q$. This is perfectly ine, but note that mathematically it's not a distance function: $\operatorname{D}\left(p\| q\right)$ is usually not equal to$\operatorname{D}\left(q\| p \right)$.

\subsection{Pinsker's Inequality}

If $p=\left(p_1,\ldots,p_m\right)$ and $q=\left(q_1,\ldots,q_m\right)$ are two distributions on $m$ elements 
\[
\| p-q\|_1 = \sum_{i=1 }^{m }|p_i-q_i|_1 \leq \sqrt{2\ln 2 \cdot \operatorname{D}\left(p\| q\right)}
\]
The other direction is not correct: $\| p-q\|_1$ may be very small and still $\operatorname{D}\left(p\| q\right)$ very large (even infinite)

\section{Entropy of a Random Variable over the Integers}

$X$: a discrete random variable (possibly infinite number of states), say over the integers 
\[
H(X) = - \sum_{i=-\infty }^{\infty} \operatorname{P}\left(X=i \right)\log \operatorname{P }\left(X=i\right) = -\sum_i \operatorname{P}_i \log \operatorname{P }_i
\]
(as in the finite case)

\subsection{Entropy of a continuous random variable}

If $X$ is a continuous random variable, we define 
\[
\operatorname{H}(X) = - \int f(x) \log f(x) \, dx
\]
where $f(x)$ is the probability density function (in analogy to the finite case). We assume that $X$ has a ``nice'' probability density function and that the integral is well defined. This is also called \textbf{Differential Entropy}. It may be negative, meaning that the equalities generalize but inequalities do not necessarily generalize.

\subsection{Intuition}

Assume that we know $X$ up to precision $\Delta$, where $\Delta$ is very small (say $\Delta=0.01$ (although the constant is meaningless)). Let $X_\Delta$ be $X$ rounded to precision $\Delta$
\[
\operatorname{H}\left(X_\Delta\right) = -\sum_{i=-\infty}^{\infty} \operatorname{P}_i \log \operatorname{P}_i
\]
where $\operatorname{P}_i=\operatorname{P}\left(i\cdot \Delta \leq X\leq \left(i+t\right)\cdot \Delta\right)$. Let $x_i$ be some point in the interval $\left[i\cdot \Delta, \left(i+1\right)\cdot \Delta\right]$. If $f$ is sufficiently nice: $\operatorname{P}_i\approx f(x_i)\cdot \Delta$, then
\begin{align*}
    \operatorname{H}\left(X_\Delta\right) &= -\sum_{i=-\infty}^{\infty} \operatorname{P}_i \log \operatorname{P}_i\\
    &\approx -\sum_{i=-\infty}^{\infty} f(x_i)\cdot \Delta\cdot \log \left(f(x_i)\cdot \Delta\right)\\
    &= -\sum_{i=-\infty}^{\infty} f(x_i)\cdot \Delta\cdot \log \left(f(x_i) + \log \Delta\right)\\
    &= -\sum_{i=-\infty}^{\infty} f(x_i)\cdot \Delta\cdot \log f(x_i)-\sum_{i=-\infty}^{\infty} f(x_i)\cdot \Delta \cdot \log \Delta\\
\end{align*}
In the limit $\Delta\to 0$, we get 
\[
\sum_{i=-\infty}^{\infty} f(x_i)\cdot \Delta\cdot \log f(x_i) \to \int f(x)\log f(x)\, dx
\]
and 
\[
\sum_{i=-\infty}^{\infty} f(x_i)\cdot \Delta \to \int f(x)\, dx = 1
\]
Therefore,
\[
\operatorname{H}\left(X_\Delta\right) + \log\Delta \to -\int f(x) \log f(x) \, dx
\]
Thus, we define 
\[
\operatorname{H}(X) = -\int f(x)\log f(x) \, dx
\]
but it should be interpreted as 
\[
\operatorname{H}(X) = -\int f(x)\log f(x) \, dx + \text{ constant}
\]
where the constant $-\log\Delta$, depends on the accuracy $\Delta$. Note that 
\[
-\log \Delta \to \infty
\]
as $\Delta \to 0$.

\subsection{Differential Entropy may be Negative}

Note that $\operatorname{H}(X)$ may be infinite, and also negative. Example $f(x)= \frac{1}{x}$ ... and I lost it.

Because differential entropy may be negative, we want to work in relative entropy.

\subsection{Relative Entropy in the Continuous Case}

If $f,g$ are sufficiently nice probability density functions 
\[
\operatorname{D}\left(f\| g\right) = \int f(x) \log \left(\frac{f(x)}{g(x)}\right)\, dx \geq 0
\]
Note that there is no uniform distribution over the real numbers.



\end{document}